\chapter{Practical 7: Master Viva Voce \& Practical Exam Revision Guide}

\section{Source Material Truncation Notice}

\begin{commonmistake}[Authoritative Source Material Truncation Notice]
The primary reference document (\texttt{ECE 279 Compressed.pdf}) contains \textbf{63 total pages}. Page 63 terminates with a navigation button \texttt{NEXT Practical 7}. The text, circuit diagrams, procedures, and viva questions for \textbf{Practical 7} were omitted in the source PDF file provided. Per strict project guidelines against AI speculation, Practical 7 is explicitly recorded here as \textbf{Truncated / Omitted in Source PDF}.

To ensure maximum utility for LPU students preparing for practical examinations, this chapter provides a comprehensive, high-yield \textbf{Master Viva Voce Revision Bank} covering all practical topics in ECE279.
\end{commonmistake}

\section{Master Viva Voce Question Bank (All Units)}

\subsection{Unit 1: Electronic Components \& Instruments (Practical 1)}

\begin{vivaqa}[Resistor Color Coding]
\textbf{Question:} State the color sequence for a $47\text{ k}\Omega \pm 10\%$ resistor.\\[4pt]
\textbf{Answer:} Yellow (4), Violet (7), Orange ($\times 10^3$), Silver ($\pm 10\%$).
\end{vivaqa}

\begin{vivaqa}[Multimeter Internal Resistance]
\textbf{Question:} What is the ideal internal resistance of a Voltmeter and an Ammeter?\\[4pt]
\textbf{Answer:} Ideal Voltmeter = $\infty$ (Infinite resistance to prevent drawing current). Ideal Ammeter = $0\,\Omega$ (Zero resistance to prevent voltage drop).
\end{vivaqa}

\begin{vivaqa}[CRO Time-Base Generator]
\textbf{Question:} What is the function of the time-base generator in a CRO?\\[4pt]
\textbf{Answer:} It generates a linear sawtooth voltage wave that sweeps the electron beam horizontally across the screen at a uniform velocity.
\end{vivaqa}

\subsection{Unit 2: DC Circuit Laws \& KVL (Practical 2)}

\begin{vivaqa}[KCL vs. KVL]
\textbf{Question:} Differentiate between Kirchhoff's Current Law (KCL) and Kirchhoff's Voltage Law (KVL).\\[4pt]
\textbf{Answer:} KCL is based on the \textbf{Law of Conservation of Charge} ($\sum I_{in} = \sum I_{out}$ at a node). KVL is based on the \textbf{Law of Conservation of Energy} ($\sum V = 0$ around a closed loop).
\end{vivaqa}

\begin{vivaqa}[Mesh vs. Loop]
\textbf{Question:} What is the difference between a loop and a mesh?\\[4pt]
\textbf{Answer:} A loop is any closed conducting path in a circuit. A mesh is a basic loop that does not contain any other loop inside it.
\end{vivaqa}

\subsection{Unit 3: Semiconductor Diodes (Practical 3)}

\begin{vivaqa}[Cut-in Voltage Factors]
\textbf{Question:} Why is cut-in voltage lower for Germanium ($0.3\text{V}$) than Silicon ($0.7\text{V}$)?\\[4pt]
\textbf{Answer:} Germanium has a smaller energy bandgap ($E_g \approx 0.67\text{ eV}$) compared to Silicon ($E_g \approx 1.12\text{ eV}$), requiring less thermal/external energy to overcome the potential barrier.
\end{vivaqa}

\begin{vivaqa}[Static vs. Dynamic Resistance]
\textbf{Question:} How does static resistance ($R_{dc}$) differ from dynamic resistance ($r_{ac}$)?\\[4pt]
\textbf{Answer:} Static resistance ($R_{dc} = V/I$) is measured at a single DC operating Q-point. Dynamic resistance ($r_{ac} = \Delta V / \Delta I$) is the reciprocal of the slope of the $V$-$I$ curve at that Q-point for AC signals.
\end{vivaqa}

\subsection{Unit 4: Digital Logic \& Universal Gates (Practical 4)}

\begin{vivaqa}[TTL 7402 IC Trap]
\textbf{Question:} Why do students often miswire IC 7402 during practical examinations?\\[4pt]
\textbf{Answer:} IC 7402 (NOR) has a reversed pin layout where Pin 1 is the \textbf{Output} and Pins 2 \& 3 are \textbf{Inputs}, unlike standard 74xx ICs (like 7400 NAND) where Pins 1 \& 2 are Inputs and Pin 3 is Output.
\end{vivaqa}

\begin{vivaqa}[Universal Gate Minimization]
\textbf{Question:} Why are universal gates preferred in industrial IC manufacturing?\\[4pt]
\textbf{Answer:} Manufacturing a single gate type (e.g., all NANDs) in mass production significantly reduces semiconductor fabrication costs, mask complexity, and inventory overhead.
\end{vivaqa}

\subsection{Unit 5: Arithmetic Adders (Practical 5)}

\begin{vivaqa}[Full Adder using 2-Input NAND Gates]
\textbf{Question:} How many 2-input NAND gates are required to construct a Full Adder?\\[4pt]
\textbf{Answer:} \textbf{9 NAND gates} (5 NAND gates per Half Adder $\times 2$, minus 1 redundant gate optimization).
\end{vivaqa}

\begin{vivaqa}[Carry Lookahead vs. Ripple Carry]
\textbf{Question:} Why is a Carry Lookahead Adder faster than a Ripple Carry Adder?\\[4pt]
\textbf{Answer:} A Ripple Carry Adder has $O(n)$ propagation delay as carry ripples sequentially through $n$ stages. A Carry Lookahead Adder calculates carry signals in parallel using $O(1)$ logic expressions ($P_i = A_i \oplus B_i, G_i = A_i B_i$).
\end{vivaqa}

\subsection{Unit 6: Sequential Flip-Flops (Practical 6)}

\begin{vivaqa}[SR to JK Transformation]
\textbf{Question:} How does a JK Flip-Flop eliminate the invalid state ($S=1, R=1$) of an SR Flip-Flop?\\[4pt]
\textbf{Answer:} By feeding back the outputs $Q$ and $\bar{Q}$ to the input steering gates, converting the invalid condition into a deterministic \textbf{Toggle state} ($Q_{n+1} = \bar{Q}_n$).
\end{vivaqa}

\begin{vivaqa}[Master-Slave Flip-Flop Operation]
\textbf{Question:} Describe how a Master-Slave JK Flip-Flop prevents Race-Around.\\[4pt]
\textbf{Answer:} It consists of two flip-flops (Master and Slave) triggered on opposite clock levels/edges. The Master responds while Clock is HIGH, but the Slave output changes only on the falling clock edge when Master is isolated, preventing feedback oscillations.
\end{vivaqa}

\section{Summary Formula \& IC Quick Reference}

\begin{table}[htbp]
\centering
\small
\begin{tabularx}{\textwidth}{l p{4.5cm} X}
\toprule
\textbf{IC Package} & \textbf{Function} & \textbf{Key Pinout / Logic Rule} \\
\midrule
\textbf{IC 7400} & Quad 2-Input NAND & Pins 1,2 In $\to$ 3 Out; Pin 7 GND, Pin 14 Vcc. \\
\textbf{IC 7402} & Quad 2-Input NOR & Pin 1 Out $\leftarrow$ Pins 2,3 In; Pin 7 GND, Pin 14 Vcc. \\
\textbf{IC 7408} & Quad 2-Input AND & Pins 1,2 In $\to$ 3 Out; Pin 7 GND, Pin 14 Vcc. \\
\textbf{IC 7410} & Triple 3-Input NAND & Gate 1: Pins 1,2,13 In $\to$ 12 Out; Gate 2: Pins 3,4,5 In $\to$ 6 Out. \\
\textbf{IC 7432} & Quad 2-Input OR & Pins 1,2 In $\to$ 3 Out; Pin 7 GND, Pin 14 Vcc. \\
\textbf{IC 7486} & Quad 2-Input XOR & Pins 1,2 In $\to$ 3 Out ($A \oplus B$); Pin 7 GND, Pin 14 Vcc. \\
\bottomrule
\end{tabularx}
\caption{Master Integrated Circuit (IC) Quick Reference Table.}
\label{tab:master-ic-ref}
\end{table}

\begin{examtip}[Practical Exam Checklist]
1. Always verify IC power connections (Pin 14 to $+5\text{V}$, Pin 7 to $\text{GND}$) before turning on power supply!\\
2. Use current limiting resistors ($220\,\Omega\text{--}470\,\Omega$) in series with all indicator LEDs to avoid burning out IC outputs.\\
3. When measuring DC voltages, connect DMM in parallel and verify probe polarity to avoid negative reading sign errors.\\
4. Double check logic level inputs: Logic $0 = 0\text{V (GND)}$, Logic $1 = +5\text{V (Vcc)}$. Never leave TTL inputs floating!
\end{examtip}
